\section{Conclusions}
\label{sec:conclusions}
\label{sec:discussion}
\label{sec:limitations}

Micrographs of strand networks show crossings that merge distinct strands and
gaps that split one strand into fragments, so each strand's identity is lost
in the image. Recovering it is what makes per-strand length, orientation and
curvature measurable. PLECTA recovers it from a binary image marking strand
centrelines alone, deciding at each junction which segments continue into one
another and allowing a strand to end rather than be forced onward. Crossing
pixels are assigned to every strand passing through them. PLECTA's grouping
core requires no training data and produces identical output from identical
input masks. The upstream segmenter is trained on synthetic image--mask pairs,
so the full workflow requires no manually annotated training micrographs.

\begin{itemize}[leftmargin=1.5em,label=\textbullet,topsep=6pt,itemsep=2pt,parsep=0pt]
  \item Across \PlectaHeldoutSceneCount{} synthetic scenes, PLECTA achieved a mean
        pairwise \Fone{} (agreement on whether two fragments belong to the
        same strand) of \PlectaHeldoutFOneMean{} [\PlectaHeldoutFOneCILow{},
        \PlectaHeldoutFOneCIHigh{}], falling from
        \PlectaDensityTwoZeroFOneMean{} to \PlectaDensitySixZeroFOneMean{} as
        areal coverage, the fraction of the frame the strands fill, rises from 20
        to 60\,\%.
  \item With every method given the same input mask, PLECTA achieved the highest
        scores among the evaluated methods across the tested densities. Its lead
        was also retained on GraFT's released generator at published settings.
  \item The ablation analysis identified the junction term penalising turning, and
        the reconnection of fragments across signal gaps, as the largest
        contributors to segmentation performance.
  \item The greyscale depth stage ordered the crossings it did not decline with
        \PlectaDepthDecidedLow{}--\PlectaDepthDecidedHigh{} accuracy, against
        chance level \PlectaDepthChance{}. Its global order is limited to
        strands connected through crossings.
\end{itemize}

\noindent\textbf{Scope and limitations.}

\begin{itemize}[leftmargin=1.5em,label=\textbullet,topsep=6pt,itemsep=2pt,parsep=0pt]
  \item PLECTA assumes directional continuity through crossings and remains
        sensitive to upstream mask errors. It can bridge short gaps but cannot
        recover omitted strands or reliably correct false axes.
  \item The real-material evaluation is a case study of
        \PlectaRealFieldCountLower{} CNT SEM fields. The automatic-axis mean
        \Fone{} of \PlectaRealUNetFOneMean{} lies within the range over which
        two human annotators agree with each other,
        \PlectaInterAnnotatorAgreeFLow{}--\PlectaInterAnnotatorAgreeFHigh{},
        so it should be interpreted alongside ambiguity in defining strand
        identity through local occlusions.
  \item The comparisons test all methods on the same binary masks in the
        dense, crossing networks studied here. They show PLECTA's advantage in
        this setting, but do not assess each method in the imaging application
        for which it was originally developed.
\end{itemize}

\noindent\textbf{Outlook.}

\begin{itemize}[leftmargin=1.5em,label=\textbullet,topsep=6pt,itemsep=2pt,parsep=0pt]
  \item \textbf{Improved binary segmentation.} A more accurate upstream segmenter,
        trained on synthetic data better matched to real micrographs, would
        improve the masks PLECTA receives.
  \item \textbf{Validation on real data.} CNT sheets stacked in a known order, or
        focused-ion-beam cross-sections of the imaged region, would give the true
        crossing order for the depth stage; a real development set, kept separate
        from the evaluation fields, would let parameters be tuned on imaging rather
        than on synthetic scenes alone.
  \item \textbf{Broader application.} Other materials, magnifications and SEM
        applications.
  \item \textbf{Mechanical models.} The segmented strands, with their depth order
        and widths, can serve as input geometry for mesoscale mechanical
        simulations.
\end{itemize}
